\documentclass[oneside,a4paper,14pt]{extarticle} %размер шрифта 14
\usepackage[T1,T2A]{fontenc}
\usepackage[
	a4paper,
	letterpaper,
	top=2cm,
	bottom=2cm,
	left=2.5cm,
	right=1.5cm
]{geometry}
\usepackage[utf8]{inputenc} %кодировка текста
\usepackage[russian]{babel} %поддержка русского языка
\usepackage{textcomp} %текстовые символы
\usepackage{indentfirst} %корректировка отступов
\usepackage{graphicx} %работа с изображениям
\usepackage{listings} %для кода !!!
\usepackage{listingsutf8}
% \usepackage{hyperref} % для ссылок
\usepackage{xcolor}
\usepackage{float} % для "впихивания" изображений в странницу
\usepackage{mwe} % for blind text and example-image-a in example
\usepackage{wrapfig}
\usepackage{caption}
\usepackage{amsmath}  % для формул и символов
\usepackage{amsfonts}
\usepackage{amsthm}
\usepackage[all]{xy}
\usepackage[breaklinks]{hyperref}
%%размер кегля у заголовков разделов
\usepackage{titlesec}
\titleformat{\section}
{\normalsize\bfseries}
{\thesection} {1em}{}
\titleformat{\subsection}
{\normalsize\bfseries}
{\thesubsection} {1em}{}
\titleformat{\subsubsection}
{\normalsize\bfseries}
{\thesubsection} {1em}{}

\renewcommand\baselinestretch{1.33}\normalsize % межстрочный интервал
\setlength{\parindent}{1.25cm}
\usepackage{indentfirst}


% прочее
\usepackage{fancyhdr}
\usepackage{colortbl}
\usepackage[dvipsnames]{xcolor}
\usepackage{tikz -timing}
\usepackage{tikz}
\usetikzlibrary{karnaugh}
\usepackage{karnaugh-map}
% \usepackage{showframe}

% настройки листинга
\lstset{
  language=C++,
  basicstyle=\ttfamily\small,
  keywordstyle=\color{blue},
  commentstyle=\color{green!60!black},
  stringstyle=\color{red},
  numbers=left,
  numberstyle=\tiny\color{gray},
  frame=single,
  breaklines=true,
  showstringspaces=false,
  tabsize=2,
  inputencoding=cp1251,
  extendedchars=true,
}

\hypersetup{
  colorlinks=true,
  linkcolor=cyan,    % цвет внутренних ссылок
  urlcolor=blue,    % цвет URL
  filecolor=magenta % цвет ссылок на файлы
}


\begin{document}
	\newpage\thispagestyle{empty}
	\begin{center}
		МИНИСТЕРСТВО НАУКИ И ВЫСШЕГО ОБРАЗОВАНИЯ\\
			РОССИЙСКОЙ ФЕДЕРАЦИИ
			ФЕДЕРАЛЬНОЕ ГОСУДАРСТВЕННОЕ БЮДЖЕТНОЕ\\
			ОБРАЗОВАТЕЛЬНОЕ
			УЧРЕЖДЕНИЕ ВЫСШЕГО ОБРАЗОВАНИЯ\\
			«ВЯТСКИЙ ГОСУДАРСТВЕННЫЙ УНИВЕРСИТЕТ»\\
			Институт математики и информационных систем\\
			Факультет автоматики и вычислительной техники\\
			Кафедра электронных вычислительных машин
	\end{center}
	\vspace{20mm}

	\begin{center}
		Отчёт по практической работе №3\\
		по дисциплине\\
		<<Введение в аппаратное обеспечение \\ систем искусственного интеллекта>>\\
		% <<Вариант 3>>\\
	\end{center}
	\vspace{48mm}

	Выполнил студент гр. ИВТб-1303-06-00 \hspace{10mm} \rule[-0,5mm]{23mm}{0.15mm}\,/Гортоломей И.К./


	Проверил доцент кафедры ЭВМ \hfill  \rule[-0,5mm]{30mm}{0.15mm}\,/Коржавина А.С./

	\vfill
	\begin{center}
		Киров\\
		2026
	\end{center}

	\newpage%\thispagestyle{empty}

	\section*{Цель работы:} Сборка схем без программирования на счётчиках, регистрах, шифраторах, дешифраторах и базовых логических элементах (НЕ, И, ИЛИ, Искл. ИЛИ ...)

	\section*{Решение заданий}

	\noindent
	\textbf{Полная таблица истинности для задания <<Детектор секретного кода>>}
\begin{center}
\begin{tabular}{cccc|cccc}
$x_1$&$x_2$&$x_3$&$x_4$&$ -x_3=q_1$&$x_1 \wedge x_2=q_2$&$q_1 \wedge x_4=q_3$&$q_2 \wedge q_3=f$\\
\hline
$0$&$0$&$0$&$0$&$1$&$0$&$0$&$0$\\
$0$&$0$&$0$&$1$&$1$&$0$&$1$&$0$\\
$0$&$0$&$1$&$0$&$0$&$0$&$0$&$0$\\
$0$&$0$&$1$&$1$&$0$&$0$&$0$&$0$\\
$0$&$1$&$0$&$0$&$1$&$0$&$0$&$0$\\
$0$&$1$&$0$&$1$&$1$&$0$&$1$&$0$\\
$0$&$1$&$1$&$0$&$0$&$0$&$0$&$0$\\
$0$&$1$&$1$&$1$&$0$&$0$&$0$&$0$\\
$1$&$0$&$0$&$0$&$1$&$0$&$0$&$0$\\
$1$&$0$&$0$&$1$&$1$&$0$&$1$&$0$\\
$1$&$0$&$1$&$0$&$0$&$0$&$0$&$0$\\
$1$&$0$&$1$&$1$&$0$&$0$&$0$&$0$\\
$1$&$1$&$0$&$0$&$1$&$1$&$0$&$0$\\
$1$&$1$&$0$&$1$&$1$&$1$&$1$&$1$\\
$1$&$1$&$1$&$0$&$0$&$1$&$0$&$0$\\
$1$&$1$&$1$&$1$&$0$&$1$&$0$&$0$\\

\end{tabular}
\end{center}

	% \newpage
	\section*{Выводы}
	Сборка <<на живую>> позволяет понять как работают базовые элементы.
	Также позволяет закрепить знания о апаратном обеспечении.

\end{document}